% !TEX root = chapter-standalone.tex

\section{Actions}
\label{sec:actions}

\todotext{Expand this intro a bit to explain the conceptual idea of actions better, and mention the key term ``state space''. Maybe move the robot on a grid example up to this intro as a concrete motivating example. Also, maybe we can motivate/explain why it makes sense to abstractly separate the set of ``instructions'' from the functions that model the dynamics... why is this a good idea? what are the benefits? maybe a second example can also underscore this, maybe one from linear algebra?}

\todotext{Because of our semicolon infix notation for function composition, the `more natural' version of actions to prefer is right actions rather than left actions (normally it is the other way around). So maybe we should start with right actions and prefer/emphasize them in our exposition (while of course still mentioning left actions and explaining them...)}

\todotext{This section is to be rewritten, with the 'covariant' (could we also say coordinate free?) version of the definition of action as the main one. Then what is currently called a right action should be 'braided' and then called a left action (but as a remark, no longer as a definition).}

\todotext{Add a remark that explains the term/concept of 'contravariant' action (there is material on this in devel at the moment. If we use the term 'contravarian morphism' somewhere, be sure we define it.)}

\todotext{Rework the following example so that it serves as the motivating introduction to this section (then no longer in an `example' environment)}

\begin{example}[Robot on a grid]
    \label{exa:robot_on_a_grid}
    We consider a robot moving on a two-dimensional unit grid, represented as $\natnumbers^2_{[0,100]}$ (\cref{fig:robot_example}).
    \begin{marginfigure}
        \centering
        \includesag{120_robotcategory}
        \caption{Robot on a grid example. }
        \label{fig:robot_example}
    \end{marginfigure}

    The state of the robot is described by its position coordinates, as $\prstel \setin \natnumbers^2_{[0,100]}$ (for convenience, consider the first element of the state to be the horizontal coordinate and the second element to be the vertical one).
    The robot can choose from a set of input actions $\prin=\{\uparrow, \downarrow,\rightarrow,\leftarrow\}$, each representing a movement of 1 unit in the corresponding direction on the grid.
    We can write the dynamics as
    \begin{equation}
        \begin{aligned}
            \prdyn\colon \prin \cartprod \prst                & \sto \prst \\
            \tupp{\uparrow, \tup{\prsteln{1},\prsteln{2}}}    & \mapsto \tupp{\prsteln{1}, \min(\prsteln{2}+1,100)}, \\
            \tupp{\downarrow, \tup{\prsteln{1},\prsteln{2}}}  & \mapsto \tupp{\prsteln{1}, \max(\prsteln{2}-1,0)}, \\
            \tupp{\rightarrow, \tup{\prsteln{1},\prsteln{2}}} & \mapsto \tupp{\min(\prsteln{1}+1,100), \prsteln{2}}, \\
            \tupp{\leftarrow, \tup{\prsteln{1},\prsteln{2}}}  & \mapsto \tupp{\max(\prsteln{1}-1,0), \prsteln{2}}, \\
        \end{aligned}
    \end{equation}
    where the ``$\min$'' and ``$\max$'' ensure that the robot stays in within the boundaries of the grid.
    The readout can be thought of as a sensor measuring the position of the robot (the state).
    Assuming a perfect sensor, we have:
    \begin{equation*}
        \defmapperiod{\prreadout}
        {\prst}
        {\sto}
        {\prout}
        {\tupp{\prsteln{1},\prsteln{2}}}
        {\tupp{\prsteln{1},\prsteln{2}}}
    \end{equation*}
    The starting position of the robot can be specified as any~$\prstart\setin \natnumbers^2_{[0,100]}$.
\end{example}
\subsection{Monoid actions}

\begin{ctdefinition}[Monoid action]
    \label{def:monoid-cov-action}
    An \maindef{action} of a \SY{monoid}~$\monoidA = \tup{\monoidAset, \mtimesof\monoidA, \idmon_\monoidA}$ on a set~\setA is a \SY{monoid homomorphism}
    \begin{equation}
        \label{eq:act-monoid}
        \act \colon \monoidA \fto \Endof\setA.
    \end{equation}
\end{ctdefinition}

Unpacked, the above definition means that we have a function
\begin{equation}
     \act \colon  \monoidAset \to (\setA \to \setA)
\end{equation}
that satisfies the following two conditions.
 
First, 
\begin{equation}
    \label{eq:act-to-endo-comp-compatibility}
    \act(\monela \mtimesof\sgrpA \monelb) = \act(\monela) \mtimesof{\Endof\setA} \act(\monelb), 
\end{equation}
which encodes ``compatibility of $\act$ with the composition operations in the source and the target'';

Second, 
\begin{equation}\label{eq:act-neutral-element-condit}
    \act(\idmon_\monoidA) = \idmon_\setA,
\end{equation}
which encodes ``compatibility of $\act$ with the neutral elements in the source and the target''. 

The latter condition may be unpacked even further; in terms of individual elements of $\setA$ it says
\begin{equation}
    \forall \setAel \setin \setA \colon \quad \act(\idmon_\monoidA)(\setAel) = \setAel.
\end{equation}

\todotext{Connect to examples now here...}

\begin{example}[robot on a grid]
\todotext{connect the intro to the definition}
\end{example}



\todotext{Discuss with AC about how to handle such ``intuitive'' examples that don't fit nicely with our composition notation/convention and our action definition convention? The same issue with be present for any matrix monoids ? If we want to keep our current conventions, maybe this is an example for motivating contravariant actions?}
\begin{example}[Scalar multiplication]
Let $\monoidA = \tup{\reals, \cdot, 1}$ be the monoid of real numbers, with multiplication as composition operation and the number $1$ as neutral element. Consider the set $\setA = \reals^2$ of real 3-dimensional vectors. An action 
\begin{equation}
\act \colon \reals \fto \Endof\reals^3
\end{equation}
is given by scalar multiplication: for each $\lambda \in \reals$, 
\begin{equation}
\defmapperiod{\act(\lambda)}{\reals^3}{\to}{\reals^3}{\begin{pmatrix} x_1 \\ x_2 \\ x_3 \end{pmatrix}}{\begin{pmatrix} x_1 \cdot \lambda \\ x_2 \cdot \lambda \\ x_3 \cdot \lambda \end{pmatrix}}
\end{equation}
\end{example}



\begin{example}[Plant growing]
\todotext{connect the example in the semigroups section that describes states of a growing plant}
\end{example}






\subsection{Uncurrying actions}

Recall that (right) currying is the operation of applying the canonical isomorphism
\begin{equation}
    \curry \colon \setC^{\setA \cartprod \setB} \to \pars{\setC^\setB}^\setA
\end{equation}
which transforms any function of the type
\begin{equation}
    \setA \cartprod \setB \to \setC
\end{equation}
into one of the type
\begin{equation}
   \setA \to (\setB \to \setC).
\end{equation}

A monoid action involves a function of the type
\begin{equation}
     \act \colon  \monoidAset \to (\setA \to \setA).
\end{equation}
If we apply the inverse $\curry^{-1}$ (so-called ``uncurrying'') to this function, 
we obtain a function of the type
\begin{equation}\label{eq:uncurried-action}
 \makecartprod{\monoidAset, \setA} \sto \setA.
 \end{equation}
Throughout mathematics and applications, actions are often presented in this uncurried form. 
Since we can tell the difference between $\act$ and its uncurried version based on their respective domains and codomains, we will use the same name, ``$\act$'', for both versions.


Now, so far, we have just used $\curry^{-1}$ to produce the function \cref{eq:uncurried-action} from a monoid action. However, $\act \colon  \monoidAset \to (\setA \to \setA)$ was not just any function, it was a monoid homomorphism $\act \colon \monoidA \fto \Endof\setA$. What does that condition translate to in the uncurried version $\act \colon  \makecartprod{\monoidAset, \setA} \sto \setA$? The answer is summarized in the following definition.

\begin{ctdefinition}[Monoid action, curried]
\label{def:uncurried-action}
   
    An \emph{uncurried action} of a \SY{monoid}~$\monoidA = \tup{\monoidAset, \mtimesof\monoidA, \idmon_\monoidA}$ on a set~\setA is:
    
    \constit
    \begin{enumerate}
        \item a function
              \begin{equation}
                  \label{eq:ract-monoid}
                  \act \colon \makecartprod{\monoidAset, \setA} \sto \setA.
              \end{equation}
    \end{enumerate}
    \condit
    \begin{enumerate}
        \item for all~$\setAel \setin \setA$ and $\monela,\monelb\setin\monoidAset$,
              \begin{equation}
                  \label{eq:monoiad-ract-1cond}
                  \act(\monela \mtimesof\monoidA \monelb, \setAel) = \act(\monelb, \act(\monela, \setAel));
              \end{equation}
        \item for all~$\setAel \setin \setA$,
              \begin{equation}
                  \label{eq:monoiad-ract-2cond}
                  \act(\idmon_\monoidA, \setAel) = \setAel.
              \end{equation}
    \end{enumerate}
\end{ctdefinition}

\begin{gradedexercise}[\exname{UncurriedAction}]
Throughout this exercise, fix a monoid $\monoidA = \tup{\monoidAset, \mtimesof\monoidA, \idmon_\monoidA}$ and~a set $\setA$. 
\begin{enumerate}
\item Given a function
    \begin{equation}
    \act \colon \makecartprod{\monoidAset, \setA} \sto \setA
    \end{equation}
    that satisfies \cref{def:uncurried-action}, prove that the curried function
    \begin{equation}
    \curry(\act) \colon  \monoidAset \to (\setA \to \setA)
    \end{equation}
    satisfies conditions \cref{eq:act-to-endo-comp-compatibility} and \cref{eq:act-neutral-element-condit} and is hence a monoid action. 
    \item Given a monoid action 
    \begin{equation}
    \act \colon  \monoidAset \to (\setA \to \setA)
    \end{equation}
    prove that its uncurried version 
    \begin{equation}
    \curry^{-1}(\act) \colon  \makecartprod{\monoidAset, \setA} \sto \setA
    \end{equation}
    satisfies \cref{def:uncurried-action}. 
\end{enumerate}
\end{gradedexercise}

\devel{

\todotext{Rewrite this part to discuss the possibility of writing uncurried action with $\monoidAset$ on the right...}

\begin{remark}
    Left actions are often denoted using infix notation of the kind
    \begin{equation}
        \makecartprod{\monoidAset, \setA} \sto \setA, \ \tup{\monela, \setAel} \mapsto \monela \cdot \setAel
    \end{equation}
    or even simply
    \begin{equation}
        \makecartprod{\monoidAset, \setA} \sto \setA, \ \tup{\monela, \setAel} \mapsto \monela \setAel
    \end{equation}
    (in regard to the latter, think of the notation often used for multiplication of matrices with vectors).
    If we write our notation ``$\monela \mtimesof\monoidA \monelb$'' as ``$\monela \star \monelb$'' (or even simply ``$\monela \monelb$''), then REF reads as
    \begin{equation}
        \monela \cdot (\monelb \cdot \setAel) = (\monela \star \monelb) \cdot \setAel,
    \end{equation}
    or simply
    \begin{equation}
        \monela (\monelb \setAel) = (\monela  \monelb)  \setAel.
    \end{equation}
    In other words, we can view REF as a kind of ``associative law'', but where two of the three elements come from $\monoidAset$, while one of the elements comes from $\setA$.
\end{remark}
}




\devel{
    \subsection{Contravariant actions}
    
    \begin{gradedexercise}[\exname{LeftActionCurry}]
    Show that having a left action of a monoid $\monoidA$ on a set $\setA$ is equivalent to having a monoid homomorphism
    \begin{equation}
        \act \colon \monoidA\op \mto \Endof\setA.
    \end{equation}
\end{gradedexercise}


\

\

    There are two versions for actions: a ``covariant'' and a ``contravariant'' version, depending on the order of the transformations.
    \footnote{
        Traditionally, people have used the names ``left'' and ``right'' to describe the two types.
        However, those names are connected to a particular choice of notation.
        We name the two versions by highlighting the properties of covariance and contravariance.
    }

    \begin{marginfigure}
        \centering
        \includesag{semigroup-contrav-action}
        \caption{Commutativity}
        \label{fig:semigroup-contrav-action}
    \end{marginfigure}
    \begin{ctdefinition}
        [Semigroup contravariant action (preliminary version)]\label{def:semigroup-contra-action-prelim}
        A \emph{contravariant semigroup action} of a \SY{semigroup}~$\sgrpA$ on a set~\setA is a map
        \begin{equation}
            \label{eq:ract}
            \ract \colon \makecartprod{\sgrpAset, \setA} \sto \setA
        \end{equation}
        such that, for all~$\setAel \setin \setA$,
        \begin{equation}
            \label{eq:ract1cond}
            \ract(\monelb, \ract(\monela, \setAel)) = \ract(\monelb \mtimesof\sgrpA \monela, \setAel).
        \end{equation}
    \end{ctdefinition}

    In contrast to \cref{eq:lact1cond}, \cref{eq:ract1cond} says that, if we first apply the action of $\monela$ to obtain $\act(\monela, \setAel)$, and then we apply the action of $\monelb$ $\act(\monelb, -)$, it is the same thing as applying the action of $\monelb\mtimes\monela$ (note the inverted order).

    % And what about contravariant actions?
    A \SY{contravariant semigroup action} can be defined as a covariant action of \emph{the opposite} \SY{semigroup}.
    And similarly for \SY{monoids} and \SY{groups}.

    \begin{ctdefinition}[Semigroup contravariant action]
        \label{def:semigroup-cont-action}
        A \maindef{contravariant semigroup action} of a \SY{semigroup}~$\sgrpA$ onto a set~\setA is a \SY{semigroup} morphism
        \begin{equation}
            \ract \colon \sgrpA\op \fto \Endof\setA.
        \end{equation}
    \end{ctdefinition}

}

\devel{
\subsection{Semigroup actions}





% There are two versions for actions: a ``covariant'' and a ``contravariant'' version, depending on the order of the transformations.
% \footnote{
%     Traditionally, people have used the names ``left'' and ``right'' to describe the two types.
%     However, those names are connected to a particular choice of notation.
%     We name the two versions by highlighting the properties of covariance and contravariance.
% }

%\linkvideo{spring2021-actions:semi-actions} % Semigroup actions

\begin{ctdefinition}[Semigroup action operation]%
    \label{def:semigroup-cov-action-prelim}
    A \emph{right action} of a \SY{semigroup}~$\sgrpA = \tup{\sgrpAset, \mthen_{\sgrpA}}$ on a set~\setA is a function
    \begin{equation}
        \label{eq:lact-semigroup}
        \act \colon \makecartprod{\setA, \sgrpAset} \sto \setA
    \end{equation}
    such that, for all~$\setAel \setin \setA$ and $\monela,\monelb\setin\sgrpAset$,
    \begin{equation}
        \label{eq:lact1cond}
        \act(\act(\setAel, \monela), \monelb) = \act(\setAel, \monela \mtimesof\sgrpA \monelb).
    \end{equation}
\end{ctdefinition}

% \begin{ctdefinition}
%     [Semigroup contravariant action (preliminary version)]\label{def:semigroup-contra-action-prelim}
%     A \emph{semigroup \textbf{contravariant} action} of a \SY{semigroup}~$\sgrpA$ on a set~$\setA$ is a map
%     \begin{equation}
%         \label{eq:ract}
%         \ract \colon \makecartprod{\sgrpAset, \setA} \to \setA
%     \end{equation}
%     such that, for all~$\setAel \setin \setA$,
%     \begin{equation}
%         \label{eq:ract1cond}
%         \ract(\monelb, \ract(\monela, \setAel)) = \ract(\monelb \mthenof\sgrpA  \monela, \setAel).
%     \end{equation}
% \end{ctdefinition}

\begin{ctdefinition}[Semigroup action]
    \label{def:semigroup-cov-action}
    An \maindef{action} of a \SY{semigroup}~$\sgrpA$ on a set~\setA is a \SY{semigroup} morphism
    \begin{equation}
        \label{eq:act-semigroup}
        \act \colon \sgrpA \fto \Endof\setA.
    \end{equation}
\end{ctdefinition}

\subsection{Actions of sets}

\begin{ctdefinition}
    [Set left action]\label{def:set-left-action-operation}
    A \emph{left action} of a set~$\sgrpAset$ on a set~\setA is a function
    \begin{equation}
        \label{eq:lact-set}
        \act \colon \makecartprod{\sgrpAset, \setA} \sto \setA.
    \end{equation}
\end{ctdefinition}

\begin{ctdefinition}%
    [Set right action]\label{def:set-right-action-operation}
    A \emph{right action} of a set~$\sgrpAset$ on a set~\setA is a function
    \begin{equation}
        \label{eq:ract-set}
        \act \colon \makecartprod{\setA, \sgrpAset} \sto \setA.
    \end{equation}
\end{ctdefinition}

\todotext{say how we can turn a set action into a semigroup action or a monoid action}

}


